\chapter{Results}

In evaluating the performance, it is important to note that the datasets used were 
specifically developed and reverse-engineered for this study to simulate real-world cases. 
A key uncertainty remains whether these synthetic datasets are sufficiently complex or 
large enough to reach the "tipping point" where the custom architecture's 
strengths truly manifest. Therefore, the evaluation should prioritize the performance 
trends across sessions as a case progresses, rather than focusing solely on 
aggregate performance metrics. \\

The fundamental trade-off addressed in this research is cost versus accuracy. 
While baseline models—which ingest the full context for every query—may offer 
higher granular precision in smaller datasets, this "stateless agent" approach 
is not scalable. In larger, multi-year cases, passing the entire history 
becomes prohibitively expensive and computationally slow. The custom architecture 
accepts a marginal risk in retrieval error to achieve sustainable 
operational costs and long-term scalability. \\

A human analogy is helpful to illustrate the architectural philosophy: 
A legal professional working on a massive case does not re-read every 
single document in the archive before answering a new question. 
Instead, they rely on a structured "factsheet" or summary (the case state) 
and perform targeted lookups when specific details are required. 
While reading every page every time might provide the most detail, 
it eventually leads to cognitive overload where one "mixes up the cards." 
By utilizing a condensed case state, the custom architecture avoids 
this "context confusion" while maintaining a manageable overhead. \\

Consequently, for the smaller datasets created for these tests, the baseline 
models often perform better because the room for retrieval error is eliminated 
when everything fits in the prompt. However, as the volume of information 
reaches realistic proportions, the necessity of the "factsheet" approach 
becomes clear. The extraction process not only manages costs but also 
yields a structured, machine-readable state that can be seamlessly 
integrated into front-end systems for human oversight, bridging the 
gap between raw data and actionable legal insights.